% =====================================================================
%  Reunión de tesis con Prof. Enrique Espina — lunes 9-11-2026, 17:00
%  Felipe Gajardo Ávalos · Magíster en Ciencias de la Ingeniería, USACH
%  Lámina 1: niveles de control — qué ENTRA, qué DECIDE, qué SALE y a quién,
%            y dónde corre cada nivel en la simulación en tiempo real
%  Lámina 2: diagrama de bloques de la operación completa (caso general,
%            Na, Nb, Nc micro-redes por fase), particionado para tiempo real
%  Lámina 3: caso simple (demanda alta prolongada en una fase)
%  Los puntos por aclarar quedan en la pauta (pauta_reunion_9nov_v3.pdf)
%  Formato: plantilla de clases USACH (fondo en fig/fondo_usach_contenido.png)
%  Compilar con XeLaTeX (dos pasadas): xelatex laminas_reunion_9nov_v3.tex
% =====================================================================
\PassOptionsToPackage{table}{xcolor}
\documentclass[aspectratio=169,10pt,t]{beamer}

\usefonttheme{professionalfonts}
\usepackage{fontspec}
\usepackage{polyglossia}
\setmainlanguage{spanish}
\usepackage{amsmath,amssymb,bm}
\usepackage{graphicx}
\usepackage{tikz}
\usetikzlibrary{arrows.meta,positioning,calc,fit,backgrounds,shapes.misc,shapes.symbols,decorations.pathreplacing}
\graphicspath{{fig/}}

% ------------------------------- Colores USACH -----------------------
\definecolor{usachT}{HTML}{009688}
\definecolor{usachTd}{HTML}{00695C}
\definecolor{usachO}{HTML}{E77500}
\definecolor{usachOl}{HTML}{FDF0E3}
\definecolor{usachR}{HTML}{C62828}
\definecolor{verdeL}{HTML}{2E7D32}
\definecolor{azulP}{HTML}{1F5FAF}
\definecolor{grisM}{HTML}{55606E}
\colorlet{cT}{usachO}\colorlet{cL}{verdeL}\colorlet{cS}{usachT}\colorlet{cP}{azulP}

\setbeamercolor{normal text}{fg=black}
\setbeamercolor{structure}{fg=usachT}
\setbeamercolor{frametitle}{fg=usachT}
\setbeamertemplate{navigation symbols}{}
\setbeamerfont{frametitle}{size=\Large}
\setbeamersize{text margin left=2.75cm,text margin right=0.55cm}
\usebackgroundtemplate{\includegraphics[width=\paperwidth,height=\paperheight]{fondo_usach_contenido.png}}
\setbeamertemplate{headline}{%
  \begin{tikzpicture}[remember picture,overlay]
    \node[fill=usachT!55,text=white,rounded corners=1.5pt,font=\tiny,inner sep=2pt,anchor=north east]
      at ([xshift=-0.22cm,yshift=-0.16cm]current page.north east) {\insertframenumber/\inserttotalframenumber};
  \end{tikzpicture}}
\setbeamertemplate{frametitle}{\vskip0.36cm\usebeamerfont{frametitle}\usebeamercolor[fg]{frametitle}\insertframetitle\par\vskip0.04cm}
\setbeamertemplate{footline}{\vbox to 0.62cm{}}

\newcommand{\SoC}{\mathrm{SoC}}

\tikzset{
  % ---- lámina 1
  tag/.style={fill=#1,text=white,rounded corners=3pt,align=center,text width=2.3cm,minimum height=1.95cm,inner sep=3pt,font=\footnotesize},
  ent/.style={draw=gray!55,fill=white,rounded corners=2pt,align=left,text width=4.55cm,minimum height=1.95cm,inner sep=4pt,font=\footnotesize,line width=.6pt},
  dec/.style={draw=#1,line width=1.3pt,fill=#1!10,rounded corners=3pt,align=center,text width=2.95cm,minimum height=1.95cm,inner sep=3pt,font=\footnotesize},
  sal/.style={draw=#1,fill=white,rounded corners=2pt,align=left,text width=4.45cm,minimum height=1.95cm,inner sep=4pt,font=\footnotesize,line width=.9pt},
  fin/.style={-{Latex[length=2.4mm]},line width=1.1pt,gray!60!black},
  fout/.style={-{Latex[length=2.4mm]},line width=1.5pt,#1},
  hdr/.style={font=\footnotesize\bfseries,text=grisM},
  rt/.style={fill=white,text=#1,rounded corners=1.5pt,font=\scriptsize\bfseries,inner sep=1.5pt},
  out/.style={draw=#1,line width=.9pt,fill=white,rounded corners=2pt,align=left,text width=2.95cm,minimum height=0.8cm,inner sep=3pt,font=\footnotesize},
  dest/.style={fill=#1,text=white,rounded corners=5pt,align=center,text width=1.7cm,minimum height=0.42cm,inner sep=2pt,font=\footnotesize\bfseries},
  fo/.style={-{Latex[length=2.2mm]},line width=1.2pt,#1},
  jd/.style={circle,fill=#1,inner sep=0pt,minimum size=3.2pt},
  bus/.style={draw=#1,fill=#1!15,line width=.9pt,rounded corners=2pt,font=\scriptsize\bfseries,text=#1!80!black,inner sep=1pt},
  halo/.style={preaction={draw=white,line width=3.6pt,-}},
  card/.style={draw=#1,line width=.8pt,rounded corners=2pt,fill=white,text width=5.65cm,inner sep=3pt,align=left,font=\fontsize{6}{7.1}\selectfont},
  ptit/.style={font=\small\bfseries,anchor=west},
  num/.style={circle,fill=#1,text=white,font=\small\bfseries,inner sep=1.2pt,minimum size=0.48cm},
  % ---- lámina 2
  band/.style={fill=#1,draw=none,rounded corners=4pt},
  bl/.style={draw=#1,line width=1.1pt,fill=white,rounded corners=2pt,align=center,font=\scriptsize,inner sep=2pt},
  pl/.style={draw=gray!60!black,line width=.7pt,fill=white,rounded corners=2pt,align=center,font=\scriptsize,inner sep=1.5pt},
  sT/.style={-{Latex[length=1.9mm]},line width=1.1pt,cT},
  sL/.style={-{Latex[length=1.9mm]},line width=1.1pt,cL,densely dashed},
  sS/.style={-{Latex[length=1.9mm]},line width=1.0pt,cS},
  sP/.style={-{Latex[length=1.9mm]},line width=1.0pt,cP},
  med/.style={-{Latex[length=1.9mm]},line width=.8pt,gray!60!black,densely dashed},
  lab/.style={font=\scriptsize},
  house/.pic={\fill[gray!70] (-0.13,0) rectangle (0.13,0.2); \fill[gray!85!black] (-0.17,0.2) -- (0,0.34) -- (0.17,0.2) -- cycle;},
}

\begin{document}

% =====================================================================
%  LÁMINA 1
% =====================================================================
\begin{frame}{Niveles de control: qué entra, qué decide y qué sale}
  \vspace{-0.28cm}
  \centering
  \resizebox{!}{6.25cm}{%
  \begin{tikzpicture}[x=1cm,y=1cm]
    % ---- encabezados
    \node[hdr] at (1.3,9.55) {NIVEL};
    \node[hdr] at (5.2,9.55) {1.\; ENTRA (mide o recibe)};
    \node[hdr] at (9.55,9.55) {2.\; DECIDE};
    \node[hdr] at (13.25,9.55) {3.\; SALE (variable)};
    \node[hdr] at (16.35,9.55) {4.\; A QUIÉN};
    \foreach \y in {7.2,5.0,2.8}{\draw[gray!30,line width=.5pt] (0,\y) -- (17.3,\y);}

    % ================= TERCIARIO (yc = 8.2) =================
    \node[tag=cT] (gT) at (1.3,8.2) {\textbf{TERCIARIO}\\[1pt] uno, central\\ cada 15 min\\[3pt] \colorbox{white}{\scriptsize\bfseries\textcolor{cT}{PC host}}};
    \node[ent] (eT) at (5.2,8.2) {%
      $\hat P_{L,\phi}$: pronóstico de clientes por fase\\
      $\hat P_{\mathrm{PV},k}$: pronóstico FV de cada MR\\
      EV$_k$: llegada, salida y $E_{\mathrm{req}}$\\
      $c(t)$: tarifa \quad $S_{h,\phi}$: cabecera\\
      \textcolor{cS}{$\SoC_k$ $\leftarrow$ secundario}};
    \node[dec=cT] (T) at (9.55,8.2) {Optimiza con ventana\\ $H=12$ h\\[2pt]
      mín.\ costo de compra\\ $+$ desbalance\\ rango de SoC en todo $t$};
    \node[out=cT] (oT1) at (13.25,8.67) {$P^*_k,\;Q^*_k$\\ {\scriptsize plan de potencia}};
    \node[out=cT] (oT2) at (13.25,7.73) {$[\SoC^{\min}_k,\;\SoC^{\max}_k]$\\ {\scriptsize rango de energía}};
    \node[dest=cP] (dT1) at (16.35,8.67) {IC$_k$};
    \node[dest=cS] (dT2) at (16.35,7.98) {Secundario $k$};
    \node[dest=cL] (dT3) at (16.35,7.45) {Lazo rápido};

    % ================= LAZO RÁPIDO (yc = 6.1) =================
    \node[tag=cL] (gL) at (1.3,6.1) {\textbf{LAZO RÁPIDO}\\[1pt] uno, central\\ cada $\sim$1 s\\[3pt] \colorbox{white}{\scriptsize\bfseries\textcolor{cL}{controlador}}};
    \node[ent] (eL) at (5.2,6.1) {%
      $i_{abc},\,v_{abc}$ medidas en la cabecera\\
      $\Rightarrow I_0$ (secuencia cero)\\
      $\Rightarrow I_2$ (secuencia negativa)\\[2pt]
      \textcolor{cT}{rango de SoC $\leftarrow$ terciario}};
    \node[dec=cL] (L) at (9.55,6.1) {Lleva $I_0,\,I_2\rightarrow 0$\\[2pt]
      corrige lo \textbf{medido},\\ no el plan\\ usa $Q$ primero, luego $P$};
    \node[out=cL] (oL1) at (13.25,6.35) {$\Delta P_k,\;\Delta Q_k$\\ {\scriptsize corrección, dentro del rango}};
    \node[dest=cP] (dL1) at (16.35,6.35) {IC$_k$};
    \node[font=\scriptsize,text=grisM,align=center] at (14.8,5.55) {reparto entre los IC: por definir};

    % ================= SECUNDARIO (yc = 3.9) =================
    \node[tag=cS] (gS) at (1.3,3.9) {\textbf{SECUNDARIO}\\[1pt] uno por MR\\ segundos\\[3pt] \colorbox{white}{\scriptsize\bfseries\textcolor{cS}{controlador}}};
    \node[ent] (eS) at (5.2,3.9) {%
      $V_{dc,k}$: tensión de su barra DC\\
      $\SoC$ de su BESS y de su EV\\
      $P_{\mathrm{PV},k}$ disponible\\[2pt]
      \textcolor{cT}{rango de SoC $\leftarrow$ terciario}};
    \node[dec=cS] (S) at (9.55,3.9) {Restaura $V_{dc,k}$\\ reparte entre sus DER\\ cuida límites de BESS y EV\\[2pt]
      \textcolor{usachR}{no ve la cabecera\\ ni otras MR}};
    \node[out=cS] (oS1) at (13.25,4.37) {$V^*_{dc,d}$\\ {\scriptsize referencia de cada DER $d$}};
    \node[out=cS] (oS2) at (13.25,3.43) {$\SoC_k$\\ {\scriptsize estado de carga de la MR}};
    \node[dest=cP] (dS1) at (16.35,4.37) {Primario DER};
    \node[dest=cT] (dS2) at (16.35,3.43) {Terciario};

    % ================= PRIMARIO (yc = 1.7) =================
    \node[tag=cP] (gP) at (1.3,1.7) {\textbf{PRIMARIO}\\[1pt] uno por convertidor\\ milisegundos\\[3pt] \colorbox{white}{\scriptsize\bfseries\textcolor{cP}{controlador}}\\[1pt] {\scriptsize potencia en sim.\ RT}};
    \node[ent] (eP) at (5.2,1.7) {%
      \textbf{DER}: $V_{dc,k}$ y \textcolor{cS}{su ref.\ $\leftarrow$ secundario}\\[2pt]
      \textbf{IC$_k$}: $P_{\mathrm{ref},k}=\textcolor{cT}{P^*_k}+\textcolor{cL}{\Delta P_k}$\\
      \phantom{\textbf{IC$_k$}: }$Q_{\mathrm{ref},k}=\textcolor{cT}{Q^*_k}+\textcolor{cL}{\Delta Q_k}$\\
      \phantom{\textbf{IC$_k$}: }$v_\phi$: tensión de su fase};
    \node[dec=cP] (P) at (9.55,1.7) {\textbf{DER}: droop $V_{dc}$--$P$\\[3pt]
      \textbf{IC$_k$}: grid-supporting\\ en AC (droop $P$--$V$\\ y $Q$--$V$ sobre $v_\phi$)};
    \node[out=cP] (oP1) at (13.25,2.17) {DER: $i_d$\\ {\scriptsize corriente hacia la barra DC}};
    \node[out=cP] (oP2) at (13.25,1.23) {IC$_k$: $p_k,\;q_k$\\ {\scriptsize potencia inyectada}};
    \node[dest=grisM] (dP1) at (16.35,2.17) {Barra DC};
    \node[dest=grisM] (dP2) at (16.35,1.23) {Fase $\phi$};
    \node[font=\scriptsize,text=grisM] at (16.35,0.7) {(la mide la cabecera)};

    % ---- 1 -> 2
    \foreach \n in {T,L,S,P}{\draw[fin] (e\n.east) -- (\n.west);}
    % ---- 2 -> 3 (abanico ortogonal con punto de bifurcación)
    \foreach \n/\c/\ya/\yb/\yc in {T/cT/8.67/7.73/8.2,S/cS/4.37/3.43/3.9,P/cP/2.17/1.23/1.7}{
      \draw[line width=1.2pt,\c] (\n.east) -- (11.4,\yc);
      \draw[fo=\c] (11.4,\yc) |- (o\n1.west);
      \draw[fo=\c] (11.4,\yc) |- (o\n2.west);
      \node[jd=\c] at (11.4,\yc) {};
    }
    \draw[fo=cL] (L.east) -- (11.4,6.1) |- (oL1.west);
    % ---- 3 -> 4
    \draw[fo=cT] (oT1.east) -- (dT1.west);
    \draw[line width=1.2pt,cT] (oT2.east) -- (15.2,7.73);
    \draw[fo=cT] (15.2,7.73) |- (dT2.west);
    \draw[fo=cT] (15.2,7.73) |- (dT3.west);
    \node[jd=cT] at (15.2,7.73) {};
    \draw[fo=cL] (oL1.east) -- (dL1.west);
    \draw[fo=cS] (oS1.east) -- (dS1.west);
    \draw[fo=cS] (oS2.east) -- (dS2.west);
    \draw[fo=cP] (oP1.east) -- (dP1.west);
    \draw[fo=cP] (oP2.east) -- (dP2.west);
  \end{tikzpicture}}

  \vspace{2pt}
  {\fontsize{6.4}{7.5}\selectfont\raggedright
  \textbf{Cómo leer:} cada fila es un nivel, de izquierda a derecha. Cada variable que sale tiene su caja y su flecha hasta el destino, con el color del nivel que la recibe. \textbf{Tiempo real:} la planta corre en el simulador; los controles, en hardware externo (I/O y red).\par}
\end{frame}

% =====================================================================
%  LÁMINA 2
% =====================================================================
\begin{frame}{Operación completa: diagrama de bloques del caso general}
  \vspace{-0.30cm}
  \centering
  \resizebox{!}{6.3cm}{%
  \begin{tikzpicture}[x=1cm,y=1cm]
    % ================= bandas de ejecución (partición para tiempo real)
    \fill[band=cT!9]    (0,8.7) rectangle (17.4,10.05);
    \fill[band=cS!8]    (0,4.6) rectangle (17.4,8.6);
    \fill[band=gray!11] (0,0)   rectangle (17.4,4.5);
    \node[rotate=90,font=\scriptsize\bfseries,cT!80!black] at (0.3,9.37) {PC host};
    \node[rotate=90,font=\scriptsize\bfseries,usachTd,align=center] at (0.3,6.6) {Controladores\\ (hardware en el lazo)};
    \node[rotate=90,font=\scriptsize\bfseries,grisM,align=center] at (0.3,2.25) {Simulador en tiempo real\\ (planta)};

    % ================= PC host
    \node[bl=cT,text width=4.6cm,minimum height=0.95cm] (T) at (9.7,9.37) {\textbf{TERCIARIO} $\cdot$ cada 15 min $\cdot$ $H=12$ h\\ plan por IC y rango de SoC por MR};
    \node[cloud,cloud puffs=11,aspect=2.6,draw=azulP,line width=.9pt,fill=white,font=\scriptsize,align=center,inner sep=0pt,minimum width=2.9cm,minimum height=1.0cm] (nu) at (14.9,9.37) {pronósticos: clientes,\\ FV, EV $\cdot$ tarifa};
    \draw[med] (nu.west) -- (T.east);

    % ================= bus de comunicación del terciario
    \node[bus=cT,minimum width=11.0cm,minimum height=0.26cm] (B) at (11.45,8.25) {bus de comunicación del terciario (cada 15 min)};
    \draw[sT] ([xshift=-0.45cm]T.south) -- ([xshift=-0.45cm]T.south |- B.north) node[pos=.5,left,lab,cT!85!black]{plan, rango};
    \draw[med] ([xshift=0.45cm]T.south |- B.north) -- ([xshift=0.45cm]T.south) node[pos=.5,right,lab,grisM]{SoC};

    % ================= alimentador 3F-4H, cabecera y clientes
    \node[draw,line width=.8pt,fill=gray!15,rounded corners=1pt,font=\scriptsize\bfseries,minimum width=0.75cm,minimum height=1.4cm] (SE) at (1.05,1.52) {S/E};
    \foreach \y/\f in {2.05/a,1.7/b,1.35/c,1.0/n}{
      \draw[line width=1.5pt] (1.43,\y) -- (17.1,\y);
      \node[font=\scriptsize\bfseries,anchor=west] at (17.1,\y) {\f};
    }
    \node[circle,draw,line width=.9pt,fill=white,font=\scriptsize\bfseries,inner sep=1pt] (M) at (2.3,1.52) {M};
    \foreach \x/\y in {3.3/2.05,4.0/1.7,4.7/1.35}{
      \fill (\x,\y) circle (1.4pt); \draw[gray!50!black,line width=.5pt] (\x,\y) -- (\x,0.62);
      \draw[gray!50!black,line width=.5pt] (\x+0.12,1.0) -- (\x+0.12,0.62); \fill (\x+0.12,1.0) circle (1.2pt);
      \pic at (\x+0.06,0.22) {house};
    }
    \node[font=\scriptsize,anchor=west,grisM] at (5.0,0.38) {clientes 1$\phi$ sin gestión (en todo el alimentador)};

    % ================= lazo rápido y mediciones de cabecera
    \node[bl=cL,text width=2.1cm,minimum height=1.0cm] (L) at (3.95,6.9) {\textbf{LAZO RÁPIDO}\\ cada $\sim$1 s\\ lleva $I_0,I_2\to0$};
    \node[lab,cL!85!black,align=center] at (3.95,5.95) {$\Delta P,\Delta Q$ a cada IC\\ (reparto por definir)};
    \draw[med] (M.north) -- (2.3,6.9) -- (L.west);
    \node[lab,anchor=east] at (2.25,6.65) {$I_0,I_2$};
    \draw[med] (2.3,6.9) -- (2.3,9.37) -- (T.west);
    \node[lab,anchor=south] at (4.6,9.37) {$S_{h,\phi}$ (cabecera)};
    \node[jd=gray!60!black] at (2.3,6.9) {};

    % ================= una columna por fase (N_phi micro-redes, copias apiladas)
    \foreach \xc/\f/\N [count=\i] in {7.65/a/N_a,11.35/b/N_b,15.05/c/N_c}{
      % --- contenedor de control de la MR
      \foreach \d in {0.16,0.08}{\draw[usachT!45,rounded corners=4pt,fill=white,line width=.6pt] (\xc-1.8+\d,4.78+\d) rectangle (\xc+1.7+\d,7.1+\d);}
      \draw[usachT!70,rounded corners=4pt,fill=white,line width=.8pt] (\xc-1.8,4.78) rectangle (\xc+1.7,7.1);
      \node[bl=cS,text width=1.5cm,minimum height=0.72cm] (S\i) at (\xc-0.55,6.65) {\textbf{Secundario}\\ MR $\f_j$};
      \node[circle,draw=cP,line width=1pt,fill=white,inner sep=0pt,minimum size=0.36cm,font=\scriptsize\bfseries] (Z\i) at (\xc+1.1,6.65) {$+$};
      \node[bl=cP,text width=1.5cm,minimum height=0.68cm] (PD\i) at (\xc-0.55,5.25) {\textbf{Primario}\\ DER (droop)};
      \node[bl=cP,text width=0.95cm,minimum height=0.68cm] (PI\i) at (\xc+1.1,5.25) {\textbf{Primario}\\ IC (g-s)};
      % --- planta: micro-red DC
      \foreach \d in {0.16,0.08}{\draw[gray!45,rounded corners=4pt,fill=white,line width=.6pt] (\xc-1.8+\d,2.55+\d) rectangle (\xc+1.7+\d,4.15+\d);}
      \draw[gray!65!black,rounded corners=4pt,fill=white,line width=.8pt] (\xc-1.8,2.55) rectangle (\xc+1.7,4.15);
      \node[pl,text width=2.05cm,minimum height=0.48cm] (DER\i) at (\xc-0.5,3.72) {PV $\cdot$ BESS $\cdot$ EV};
      \draw[line width=2pt] (\xc-1.5,3.15) -- (\xc+0.62,3.15);
      \draw[line width=1pt] (DER\i.south) -- (\xc-0.5,3.15);
      \node[pl,minimum width=0.8cm,minimum height=0.48cm] (IC\i) at (\xc+1.1,3.15) {\textbf{IC}};
      \draw[line width=1pt] (\xc+0.62,3.15) -- (IC\i.west);
      \node[font=\scriptsize\bfseries,anchor=south west] at (\xc-1.75,2.58) {MR $\f_j$, \;$j\le \N$};
      % --- secundario -> primario DER; sumador -> primario IC
      \draw[sS] (S\i.south) -- (PD\i.north);
      \draw[sP] (Z\i.south) -- (PI\i.north);
      % --- I/O entre controladores y simulador
      \draw[sP] ([xshift=0.35cm]PD\i.south) -- ([xshift=0.35cm]PD\i.south |- DER\i.north);
      \draw[med] ([xshift=-0.35cm]PD\i.south |- DER\i.north) -- ([xshift=-0.35cm]PD\i.south);
      \draw[sP] ([xshift=-0.2cm]PI\i.south) -- ([xshift=-0.2cm]IC\i.north);
      \draw[med] ([xshift=0.2cm]IC\i.north) -- ([xshift=0.2cm]PI\i.south);
      \draw[med] (DER\i.west) -- (\xc-1.68,3.72) -- (\xc-1.68,6.65) -- (S\i.west);
    }
    % --- lazo rápido -> sumador de cada IC (línea con bifurcaciones)
    \draw[line width=1.3pt,cL,densely dashed] (L.east) -- (5.5,6.9) -- (5.5,7.55) -- (15.5,7.55);
    \foreach \xc [count=\i] in {7.65,11.35,15.05}{
      \draw[sL] (\xc+0.45,7.55) -- (\xc+0.45,6.65) -- (Z\i.west);
      \node[jd=cL] at (\xc+0.45,7.55) {};
    }
    % --- terciario <-> MR por el bus (cruzan la línea verde con halo)
    \foreach \xc [count=\i] in {7.65,11.35,15.05}{
      \draw[sT,halo] (\xc-0.3,8.12) -- (\xc-0.3,7.01);
      \draw[med,halo] (\xc-0.85,7.01) -- (\xc-0.85,8.12);
      \draw[sT,halo] (\xc+1.1,8.12) -- (Z\i.north);
    }
    % --- IC -> su fase
    \draw[line width=1.3pt] (IC1.south) -- (8.75,2.05); \fill (8.75,2.05) circle (1.8pt);
    \draw[line width=1.3pt] (IC2.south) -- (12.45,1.7); \fill (12.45,1.7) circle (1.8pt);
    \draw[line width=1.3pt] (IC3.south) -- (16.15,1.35); \fill (16.15,1.35) circle (1.8pt);

    % --- rótulos una sola vez (columna a)
    \node[lab,grisM,anchor=east] at (6.75,7.82) {SoC};
    \node[lab,cT!85!black,anchor=west] at (7.4,7.82) {rango};
    \node[lab,cT!85!black,anchor=west] at (8.8,7.82) {$P^*\!,Q^*$};
    \node[lab,cS!80!black,anchor=east] at (7.05,5.95) {refs.};
    \node[lab,cP,anchor=east] at (8.7,5.95) {$P_{\mathrm{ref}},Q_{\mathrm{ref}}$};
    \node[lab,grisM,anchor=east] at (5.92,5.0) {$V_{dc},\SoC$};
    \node[lab,cP,anchor=west] at (7.5,4.4) {PWM};
    \node[lab,grisM,anchor=east] at (6.75,4.4) {$v,i$};
  \end{tikzpicture}}

  \vspace{1pt}
  {\fontsize{6.4}{7.5}\selectfont\raggedright
  \textcolor{cT}{\rule[0.5ex]{0.4cm}{1.3pt}} plan y rango (15 min)\quad
  \textcolor{cL}{\rule[0.5ex]{0.15cm}{1.3pt}\,\rule[0.5ex]{0.15cm}{1.3pt}} corrección del lazo rápido (s)\quad
  \textcolor{cS}{\rule[0.5ex]{0.4cm}{1.3pt}} secundario $\rightarrow$ DER\quad
  \textcolor{cP}{\rule[0.5ex]{0.4cm}{1.3pt}} $P_{\mathrm{ref}}=P^*+\Delta P$ y PWM\quad
  \textcolor{gray!60!black}{\rule[0.5ex]{0.15cm}{0.9pt}\,\rule[0.5ex]{0.15cm}{0.9pt}} mediciones\quad
  g-s: grid-supporting. Cada flecha termina en el bloque que la recibe; los puntos marcan bifurcaciones; copias apiladas $=N_\phi$ micro-redes por fase.\par}
\end{frame}

% =====================================================================
%  LÁMINA 3 — CASO SIMPLE, versión visual (la energía es un área)
% =====================================================================
\begin{frame}{Caso simple: ¿cuánto se corrige si las baterías no alcanzan?}
  \vspace{-0.30cm}
  \centering
  \resizebox{!}{6.55cm}{%
  \begin{tikzpicture}[x=1cm,y=1cm,font=\small]
    \colorlet{cRes}{usachR}\colorlet{cDes}{usachT}\colorlet{cCar}{verdeL}

    % ===================== 1. EL PROBLEMA =====================
    \node[num=usachO] at (0.25,9.15) {1};
    \node[ptit] at (0.55,9.15) {El problema};
    \foreach \x/\f in {0.6/b,1.4/c}{
      \fill[gray!35] (\x,6.75) rectangle (\x+0.55,7.95);
      \node[font=\small\bfseries] at (\x+0.275,6.5) {\f};}
    \fill[gray!35] (2.2,6.75) rectangle (2.75,7.95);
    \fill[cRes!75] (2.2,7.95) rectangle (2.75,8.85);
    \node[font=\small\bfseries] at (2.475,6.5) {a};
    \draw[decorate,decoration={brace,amplitude=4pt},line width=.8pt] (2.85,8.85) -- (2.85,7.95) node[midway,right=4pt,font=\small,align=left]{$D=6$ kW\\ extra};
    \draw[-{Latex[length=2.4mm]},line width=1.2pt,cRes] (4.15,7.85) -- (4.75,7.85);
    \node[draw=cRes,line width=1pt,rounded corners=3pt,fill=white,align=center,font=\small] at (5.75,7.85) {sale por\\ el neutro:\\[2pt] $I_n=\dfrac{6000}{220}$\\[3pt] $\mathbf{27{,}3\ A}$};

    % ===================== 2. EL PRESUPUESTO DE ENERGÍA =====================
    \node[num=usachO] at (7.25,9.15) {2};
    \node[ptit] at (7.55,9.15) {La energía no alcanza};
    \fill[cRes!22] (7.4,7.0) rectangle (12.2,8.5);
    \fill[cDes!70] (7.4,7.0) rectangle (9.0,8.5);
    \draw[line width=.9pt] (7.4,7.0) rectangle (12.2,8.5);
    \node[font=\small\bfseries,white,align=center] at (8.2,7.75) {8 kWh\\ \footnotesize batería};
    \node[font=\small,align=center] at (10.6,7.75) {16 kWh\\ \footnotesize sin cubrir};
    \draw[{Latex[length=2mm]}-{Latex[length=2mm]},line width=.7pt] (7.4,6.8) -- node[below,font=\footnotesize]{$T=4$ h} (12.2,6.8);
    \draw[{Latex[length=2mm]}-{Latex[length=2mm]},line width=.7pt] (7.25,7.0) -- (7.25,8.5);
    \node[font=\footnotesize,anchor=west] at (7.4,8.72) {corregir todo exige $D\cdot T=24$ kWh};
    \node[draw=cDes,line width=1pt,rounded corners=3pt,fill=white,align=center,font=\small,text width=3.6cm] at (15.9,7.75) {La batería cubre \textbf{1/3}\\[2pt] el promedio baja de\\ 27,3 A a $\mathbf{18{,}2\ A}$,\\[2pt] \textbf{haga lo que haga}};
    \draw[-{Latex[length=2.4mm]},line width=1.2pt,cDes] (12.35,7.75) -- (14.0,7.75);

    % ===================== 3. TRES FORMAS DE GASTAR LOS MISMOS 8 kWh =====================
    \node[num=usachO] at (0.25,6.0) {3};
    \node[ptit] at (0.55,6.0) {Los mismos 8 kWh (misma área), gastados de tres formas};
    % leyenda
    \fill[cRes!55] (10.7,5.88) rectangle (11.0,6.12); \node[font=\footnotesize,anchor=west] at (11.05,6.0) {queda en el neutro};
    \fill[cDes!75] (13.65,5.88) rectangle (13.95,6.12); \node[font=\footnotesize,anchor=west] at (14.0,6.0) {descarga en $a$};
    \fill[cCar!70] (15.9,5.88) rectangle (16.2,6.12); \node[font=\footnotesize,anchor=west] at (16.25,6.0) {carga en $b,c$};

    % --- macro de ejes: \px = origen x del panel
    \foreach \px/\L/\tit in {0.6/A/{gastar todo al inicio},6.6/B/{repartir parejo},12.6/C/{B + cargar en $b$ y $c$}}{
      \node[font=\small\bfseries,anchor=west] at (\px-0.4,5.35) {\L.\; \tit};
      \draw[gray!60,line width=.6pt] (\px,2.45) -- (\px+4.85,2.45);
      \draw[gray!60,line width=.6pt] (\px,2.45) -- (\px,5.05);
      \foreach \t/\l in {0/18,2.4/20,4.8/22}{\draw[gray!60] (\px+\t,2.45) -- (\px+\t,2.35) node[below,font=\footnotesize,gray!60!black]{\l:00};}
      \foreach \k in {2,4,6}{\draw[gray!60] (\px,2.45+0.4*\k) -- (\px-0.07,2.45+0.4*\k);}
      \draw[gray!65!black,dashed,line width=.7pt] (\px,4.85) -- (\px+4.8,4.85);
    }
    \foreach \k in {0,2,4,6}{\node[font=\footnotesize,gray!60!black,anchor=east] at (0.53,2.45+0.4*\k) {\k};}
    \node[font=\footnotesize,gray!60!black,rotate=90] at (-0.35,3.65) {desbalance [kW]};

    % --- A: 6 kW durante 1,33 h (1,6 cm), luego nada
    \fill[cDes!75] (0.6,2.45) rectangle (2.2,4.85);
    \fill[cRes!55] (2.2,2.45) rectangle (5.4,4.85);
    \node[font=\footnotesize\bfseries,white,align=center] at (1.4,3.65) {8 kWh\\ en 1,3 h};
    \node[font=\footnotesize\bfseries,align=center] at (3.8,3.65) {27,3 A\\ desde 19:20};
    % --- B: 2 kW parejo durante 4 h
    \fill[cRes!55] (6.6,2.45) rectangle (11.4,4.05);
    \fill[cDes!75] (6.6,4.05) rectangle (11.4,4.85);
    \node[font=\footnotesize\bfseries,white] at (9.0,4.45) {8 kWh repartidos parejo};
    \node[font=\footnotesize\bfseries] at (9.0,3.25) {18,2 A todo el evento};
    % --- C: B + carga 2 kW en b y c
    \fill[cRes!55] (12.6,2.45) rectangle (17.4,3.25);
    \fill[cCar!70] (12.6,3.25) rectangle (17.4,4.05);
    \fill[cDes!75] (12.6,4.05) rectangle (17.4,4.85);
    \node[font=\footnotesize\bfseries,white] at (15.0,4.45) {8 kWh descarga en $a$};
    \node[font=\footnotesize\bfseries,white] at (15.0,3.65) {$+$ carga 2 kW en $b$ y $c$};
    \node[font=\footnotesize\bfseries] at (15.0,2.85) {9,1 A};

    % --- veredictos
    \node[draw=usachR,line width=.9pt,rounded corners=3pt,fill=white,align=center,font=\footnotesize,text width=4.6cm] at (3.0,1.35) {pico \textbf{27,3 A} $\cdot$ RMS 22,3 A\\ \textcolor{usachR}{\textbf{gasta la energía antes de tiempo}}};
    \node[draw=usachT,line width=1.4pt,rounded corners=3pt,fill=usachT!8,align=center,font=\footnotesize,text width=4.6cm] at (9.0,1.35) {pico \textbf{18,2 A} $=$ promedio\\ \textcolor{usachTd}{\textbf{óptimo con la misma energía}}};
    \node[draw=verdeL,line width=.9pt,rounded corners=3pt,fill=white,align=center,font=\footnotesize,text width=4.6cm] at (15.0,1.35) {pico \textbf{9,1 A}\\ \textcolor{verdeL}{\textbf{mejor, pero importa 16 kWh más que B}}};

    % ===================== CONCLUSIÓN =====================
    \node[draw=usachO,line width=1.1pt,rounded corners=3pt,fill=usachOl,align=center,font=\small,text width=17.6cm,inner sep=4pt] at (8.95,0.15) {%
      \textbf{El área} (energía del rango) fija \textbf{cuánto} se corrige; \textbf{la forma} en el tiempo fija \textbf{el pico}.
      Para elegir B, el terciario debe ver el evento completo: $H\ge T$ $\Rightarrow$ $\mathbf{H=12}$ \textbf{h}. Con $H=1$ h no lo ve y termina en A.};
  \end{tikzpicture}}

  \vspace{0pt}
  {\fontsize{5.4}{6.4}\selectfont\color{grisM}\raggedright
  Supuestos: cabecera rígida ($V=220$ V), sin pérdidas, $D$ constante, $Q=0$, sin EV (con EV la batería aporta solo el margen que deja $E_{\mathrm{req}}$). En C, $b$ y $c$ necesitan 8 kWh de margen de carga cada una.\par}
\end{frame}

\end{document}
